% Source-derived chapter 00: Abstract and front matter
% Original source: intersection-complexes-unramified-l-factors
\chapter{Abstract and front matter}
\label{intersection-complexes-unramified-l-factors-chapter-00}
\source{intersection-complexes-unramified-l-factors}

\begin{remark}[Source section]
\label{intersection-complexes-unramified-l-factors-chapter-00-source}
\source{intersection-complexes-unramified-l-factors}
\source{intersection-complexes-unramified-l-factors}
This chapter reproduces the corresponding section of the retrieved
LaTeX source, including its definitions, statements, examples, and
proofs. It has not yet been translated into Lean.
\notready
\end{remark}

% The source section title was: Abstract and front matter
\begin{abstract}
% arXiv abstract
%Let X be an affine spherical variety, possibly singular, and $L^+X$ its arc space. The intersection complex of $L^+X$, or rather of its finite-dimensional formal models, is conjectured to be related to special values of local unramified L-functions. Such relationships were previously established in Braverman-Finkelberg-Gaitsgory-Mirkovic for the affine closure of the quotient of a reductive group by the unipotent radical of a parabolic, and in Bouthier-Ngo-Sakellaridis for toric varieties and L-monoids. In this paper, we compute this intersection complex for the large class of those spherical G-varieties whose dual group is equal to the Langlands dual group of G, and the stalks of its nearby cycles on the horospherical degeneration of X. We formulate the answer in terms of a Kashiwara crystal, which conjecturally corresponds to a finite-dimensional representation of the dual group determined by the set of B-invariant valuations on X. We prove the latter conjecture in many cases. Under the sheaf--function dictionary, our calculations give a formula for the Plancherel density of the IC function of $L^+X$ as a ratio of local L-values for a large class of spherical varieties.
%
Let $X$ be an affine spherical variety, possibly singular, and $\msf L^+X$ its arc space. The intersection complex of $\msf L^+X$, or rather of its finite-dimensional formal models, is conjectured to be related to special values of local unramified $L$-functions. Such relationships were previously established in \cite{BFGM} for the affine closure of the quotient of a reductive group by the unipotent radical of a parabolic, and in \cite{BNS} for toric varieties and $L$-monoids. In this paper, we compute this intersection complex for the large class of those spherical $G$-varieties whose dual group is equal to $\check G$, and the stalks of its nearby cycles on the horospherical degeneration of $X$. We formulate the answer in terms of a Kashiwara crystal, which conjecturally corresponds to a finite-dimensional $\check G$-representation determined by the set of $B$-invariant valuations on $X$. 
We prove the latter conjecture in many cases. Under the sheaf--function dictionary,
our calculations give a formula for the Plancherel density of the IC function
of $\msf L^+X$ as a ratio of local $L$-values for a large class of spherical varieties.
\end{abstract}
